\chapter{Herkömmliche Regelung der linearen und nichtlinearen Regelstrecke} \label{Kapitel_Regelung_lin_nonlin}
In diesen Kapitel wird die Auslegung der Regler für die lineare und nichtlineare \\ Regelstrecke beschrieben.
Aufgrund der Nichtlinearität des Hochsetzstellers wie in\\  Abschnitt \ref{Section_Nichtlineare_Regelstrecke} beschrieben, wird in der Auslegung kaskadierte Regelungen verwendet.
Die Kaskadenregelung besteht aus einer unterlagerten Stromregelung und einer\\ überlagerten Spannungsregelung.
Für eine bessere Vergleichbarkeit zwischen der\\ linearen und nichtlinearen Regelstrecke wird auch für die lineare Regelstrecke eine \\kaskadierte Regelung verwendet.
Ein Antiwindup kommt zum Einsatz, um den \\Integrator abzustellen, wenn die Stellreserve ausgeschöpft ist.
\\Das Einschwingverhalten der Kaskadenregelung von Tief- und Hochsetzstellern basiert auf dem Einschwingverhalten eines Tiefpass kritischer Dämpfung.
In den folgenden\\  Abschnitten beschreibt diese Arbeit die Auslegung der kaskadierten Regelung des \\Hoch- und Tiefsetzstellers.
\newpage
\section{Lineare Regelstrecke}
Für die Auslegung der kaskadierten Regelung des Tiefsetzstellers wird das \\ regelungstechnische Blockschaltbild, hergeleitet aus dem physikalischen Modell des Tiefsetzstellers (siehe Abschnitt \ref{Section_lineare_Regelstrecke}), verwendet. \\
\begin{figure}[h]
	\centering
	\includegraphics[width=12cm]{TSS_RBS.pdf}
	\caption{Regelungstechnisches Blockschaltbild des Tiefsetzstellers}
\end{figure} 

\newpage
\subsection{Auslegung der unterlagerten Stromregelung}\label{Section_Unterlagerte_Stromregelung_TSS}
Wie in Abbildung \ref{TSS_Stromregelung} dargestellt, entsteht ein Regelsystem nach dem Prinzip der \\Rückführung eines vollständigen Satzes von Zustandsvariablen (siehe Abschnitt \ref{Rückführung_vollständiger_Zustandsvariablen}).
\newline
\begin{figure}[h]
	\centering
	\includegraphics[width=16cm]{TSS_Stromregelung.pdf}
	\caption{Stromregelung des Tiefsetzstellers}
	\label{TSS_Stromregelung}
\end{figure} \\ \noindent
Die Auslegung der Einstellparameter $K$ und $K_1$ erfolgt mithilfe der\\ Übertragungsfunktion.
\begin{align}
    \ILEAvar{F}{R}(p) &= \frac{\ILEAvar{i}{L}}{\ILEAvar{i}{L,soll}}  \\
    \ILEAvar{F}{R}(p) &= \frac{K}{\ILEAvar{K}{1} K + p \ILEAvar{T}{L}} \\
    \ILEAvar{F}{R}(p) &\widehat{=} \ILEAvar{K}{w} \cdot \frac{1}{1 + p \ILEAvar{T}{Strom}} 
\end{align}
\noindent
Durch Koeffizientenvergleich und der Bedingung des Eingeschaftsparameter $\ILEAvar{K}{w} = 1$ (siehe Abschnitt \ref{Rückführung_vollständiger_Zustandsvariablen}) ergeben sich die Einstzellparameter.
\begin{align}
    \ILEAvar{K}{1} &= \frac{1}{\ILEAvar{K}{w}} \\
    K &= \frac{T_L}{\ILEAvar{K}{1} \ILEAvar{T}{Strom}} = \frac{\ILEAvar{T}{L}}{\ILEAvar{T}{Strom}}
\end{align}
Die Rückkopplung der Ausgangsspannung $\ILEAvar{u}{C}$ über den Parameter $\ILEAvar{K}{2}$ dient der \\ Auflösung der Kreisschaltung, verursacht durch die Vorkopplung der \\ Ausgangsspannung. 
Für die Störgrößenkompensation gilt.
\begin{align}
	\ILEAvar{K}{2} = -\frac{1}{K} = - \frac{\ILEAvar{T}{Strom}}{\ILEAvar{T}{L}}
\end{align}
Die unterlagerte Stromregelung muss deutlich schneller ausgelegt werden als die \\ überlagerte Spannungsregelung, damit der unterlagerte Regler als eingeschwungen betrachtet werden kann. 
Typischerweise wird ein Verhältnis der Zeitkonstanten von $\ILEAvar{T}{Strom} : \ILEAvar{T}{Spannung} \approx 1:10$ angestrebt.
\newpage
\subsection{Auslegung der überlagerten Spannungsregelung}
Aufgrund der genannten Bedingung der Zeitkonstanten in Abschnitt \ref{Section_Unterlagerte_Stromregelung_TSS} gilt die \\ Annahme, dass die unterlagerte Stromregelung bereits eingeschwungen ist.
Damit kann die Auslegung der Spannungsregelung die unterlagerte Stromregelung \\unberücksichtigt lassen und es gilt $\ILEAvar{i}{L,soll} = \ILEAvar{i}{L}$.
Der Laststrom $\ILEAvar{i}{R}$ stellt eine Störgröße dar und bleibt daher in der Berechnung der Übertragungsfunktion unberücksichtigt.
Zuerst erfolgt die Betrachtung des Systems ohne Bypass-I-Regler.
\newline
\begin{figure}[h]
	\centering
	\includegraphics[width=16cm]{TSS_Spannungsregelung_ohne_Bypass.pdf}
	\caption{Überlagerte Spannungsregelung des Tiefsetzstellers ohne Bypass-I-Regler}
	\label{Überlagerte_Spannungsregelung_ohne_Bypass_TSS}
\end{figure} \\ \noindent
Die Übertragungsfunktion des Systems ohne Bypass-I-Regler ist.
\begin{align}
	\ILEAvar{F}{R}(p)  &= \frac{u_{C}}{\ILEAvar{u}{C,soll}^{*}}  \\
	\ILEAvar{F}{R}(p) &= \frac{C}{C \ILEAvar{C}{1} + p \ILEAvar{T}{C}} \\
	\ILEAvar{F}{R}(p) &\widehat{=}\ILEAvar{C}{w}^{*} \cdot \frac{1}{1 + p \ILEAvar{T}{Sp}^{*}} 
\end{align}
Damit ergeben sich die Einstellparameter zu.
\begin{align}
	\ILEAvar{C}{1} &= \frac{C}{C \ILEAvar{C}{w}^{*}} \\
	C &= \frac{\ILEAvar{T}{C}}{\ILEAvar{T}{Sp}^{*} \ILEAvar{C}{1}}
\end{align}
\newpage \noindent
Der Bypass-I-Regler dient zur Kompensation einer stationären Sollwertabweichung und ist in Abbildung \ref{TSS_Spannungsregelung} zur Spannungsregelung hinzugefügt.
\begin{figure}[h]
	\centering
	\includegraphics[width=16cm]{TSS_Spannungsregelung.pdf}
	\caption{Spannungsregelung des Tiefsetzstellers}
	\label{TSS_Spannungsregelung}
\end{figure}  \\ \noindent
Hierzu wird die Übertragungsfunktion mit folgenden Ansätzen betrachtet wie in\\  Abschnitt \ref{Bypass} beschrieben.
\begin{align}
	\frac{1}{\ILEAvar{C}{w}^{*}} &= 1 + \frac{T_{Sp}}{T_{I}} \\
	\frac{1}{\ILEAvar{C}{w}^{*}} \ILEAvar{T}{Sp}^{*} &= T \cdot \left( 1 + \ILEAvar{q}{2} \frac{\ILEAvar{T}{Sp}}{\ILEAvar{T}{I}} \right) 
\end{align}
Damit ergibt sich die Einstellparameter.
\begin{align}
	\ILEAvar{C}{w}^{*} &= \frac{1}{1 + \frac{T_{Sp}}{T_{I}}} \\
	\ILEAvar{C}{1} &= \frac{1}{\ILEAvar{C}{w}^{*}} \\
	C &= \frac{\ILEAvar{T}{C}}{\ILEAvar{T}{Sp}^{*} \ILEAvar{C}{1}}
\end{align}
Und der Eigenschaftsparameter.
\begin{align}
	\ILEAvar{T}{Sp}^{*} &= \ILEAvar{T}{Sp} \cdot \left( 1 + \ILEAvar{q}{2} \frac{\ILEAvar{T}{Sp}}{\ILEAvar{T}{I}} \right) 
\end{align}
\newpage
\section{Nichtlineare Regelstrecke}
Für die Auslegung der kaskadierten Regelung des Hochsetzstellers dient das \\ regelungstechnische Blockschaltbild, das auf Basis des physikalischen Modells des Hochsetzstellers abgeleitet wurde (siehe Abschnitt \ref{Section_Nichtlineare_Regelstrecke}).
Die ursprünglich nichtlineare Übertragungsfunktion des Hochsetzstellers wird mittels geeigneter Umrechnungswerke linearisiert. 
Diese Umrechnungswerke erfassen die nichtlinearen Zusammenhänge \\ zwischen den relevanten Größen des Hochsetzstellers, wie in Abschnitt \ref{PhyM_HSS} \\ beschrieben, sodass im Rahmen der Reglerauslegung die Nichtlinearitäten \\ vernachlässigt werden können.
\newline
\begin{figure}[h]
	\centering
	\includegraphics[width=12cm]{HSS_RBS_2.pdf}
	\caption{Regelungstechnisches Blockschaltbild des Hochsetzstellers}
	\label{RBS_HSS2}
\end{figure} 
\newpage
\subsection{Auslegung der unterlagerten Stromregelung} \label{Section_Unterlagerte_Stromregelung}
Wie in Abbildung \ref{HSS_Stromregelung} dargestellt, entsteht ein Regelsystem nach dem Prinzip der\\  Rückführung eines vollständigen Satzes von Zustandsvariablen (siehe Abschnitt \ref{Rückführung_vollständiger_Zustandsvariablen}). \\ 
\newline
\begin{figure}[h!]
	\centering
	\includegraphics[width=16cm]{HSS_Stromregler.pdf}
	\caption{Schaltbild eines Hochsetzstellers}
	\label{HSS_Stromregelung}
\end{figure} \\ \noindent
Die Auslegung der Einstellparameter $K$ und $\ILEAvar{K}{1}$ erfolgt mithilfe der\\  Übertragungsfunktion.
\begin{align}
    \ILEAvar{F}{R}(p) &= \frac{\ILEAvar{i}{L}}{\ILEAvar{i}{L,soll}}  \\
    \ILEAvar{F}{R}(p) &= \frac{K}{\ILEAvar{K}{1}K + p\ILEAvar{T}{L}} \\
    \ILEAvar{F}{R}(p) &\widehat{=} \ILEAvar{K}{w} \cdot \frac{1}{1 + p \ILEAvar{T}{Strom}}
\end{align}
Im eingeschwungen Zustand soll $\ILEAvar{i}{L,soll} $ gleich $\ILEAvar{i}{L}$ sein (siehe \ref{Rückführung_vollständiger_Zustandsvariablen}). Daher wird $\ILEAvar{K}{w}$ wie folgt gewählt.
\begin{align}
    \ILEAvar{K}{w} &= 1 
\end{align}
\newpage \noindent
Es ergibt sich für die Einstellparameter.
\begin{align}
    \ILEAvar{K}{1} &= \frac{1}{\ILEAvar{K}{w}} \\
    K &= \frac{\ILEAvar{T}{L}}{\ILEAvar{K}{1} \ILEAvar{T}{Strom}} 
\end{align} \\ \noindent
Damit die Sollvorgabe des Stromreglers in die Stellgröße umgerechnet werden kann, muss das Steuergesetz mittels eines Umrechenwerkes beachtet werden.
\begin{figure}[h]
	\centering
	\includegraphics[width=10cm]{HSS_U2.pdf}
	\caption{Umrechenwerk}
	\label{Umrechenwerk2}
\end{figure}\\ \noindent
Der Sollwert der Spulenspannung $\ILEAvar{u}{L,soll}$ wird mit der Eingangsspannung $\ILEAvar{u}{in}$ und der Kondensatorspannung $\ILEAvar{u}{C}$ verrechnet, um die Stellgröße, die relative Einschaltdauer, $\ILEAvar{\tau}{g}$ zu erhalten.
\newpage
\subsection{Auslegung der überlagerten Spannungsregelung} \label{Section_Überlagerte_Spannungsregelung}
In Abbildung \ref{HSS_Spannungsregelung} ist die überlagerte Spannungsregelung, das Umrechenwerk und der Bypass-I-Regler dargestellt.
Die Auslegung erfolgt analog zu der Auslegung der\\ überlagerten Spannungsregelung des Tiefsetzstellers (Abschnitt \ref{TSS_Spannungsregelung}) und wird daher verkürzt dargestellt.
\\ Das Umrechenwerk berücksichtigt die aus dem Abschnitt \ref{PhyM_HSS} gezeigte Nichtlinearität zwischen $\ILEAvar{i}{L}$ und $\ILEAvar{i}{C}$.
\newline
\begin{figure}[h]
	\centering
	\includegraphics[width=16cm]{HSS_Spannungsregelung.pdf}
	\caption{Schaltbild eines Hochsetzstellers}
	\label{HSS_Spannungsregelung}
\end{figure}  \\ \noindent
Hierbei wird der Laststrom $\ILEAvar{i}{R}$ als Störgröße betrachtet und daher in der Auslegung vernachlässigt.
\begin{align}
	\ILEAvar{C}{w}^{*} &= \frac{1}{1 + \frac{\ILEAvar{T}{Sp}}{\ILEAvar{T}{I}}} \\
	\ILEAvar{T}{Sp}^{*} &= \ILEAvar{T}{Sp} \cdot \left( 1 + \ILEAvar{q}{2} \frac{\ILEAvar{T}{Sp}}{\ILEAvar{T}{I}} \right) \\
	\ILEAvar{C}{1} &= \frac{1}{\ILEAvar{C}{w}^{*}} \\
	C &= \frac{\ILEAvar{T}{C}}{\ILEAvar{T}{Sp}^{*} \ILEAvar{C}{1}}
\end{align} \newpage \noindent
Die Sollwertvorgabe des Kondensatorstroms $\ILEAvar{i}{C,soll}$, wie in Abbildung \ref{Umrechenwerk1} beschrieben, wird über das Umrechenwerk in die Sollwertvorgabe des Spulenstroms $\ILEAvar{i}{L,soll}$ \\ umgerechnet.
\begin{figure}[h]
    \centering
    \includegraphics[width=8cm]{HSS_U1.pdf}
    \caption{Umrechenwerk}
    \label{Umrechenwerk1}
\end{figure}

